\documentclass{mathnotes}

\input{course}
\notedate{September 10, 2026}

\begin{document}

\notetitle{12}{Uniform Continuity}

Let $f: A \to \R$ be a function on a set $A \subseteq \R$. Recall that continuity
at a point $c$ allows the choice of $\delta$ to depend on both $\eps$ and $c$.
Uniform continuity strengthens this by requiring a single $\delta$ that works
across the entire domain.

\begin{definition}
A function $f: A \to \R$ is uniformly continuous on $A$ if for every
$\eps > 0$ there exists a $\delta > 0$ such that for all $x, y \in A$,
\[
  |x - y| < \delta \implies |f(x) - f(y)| < \eps.
\]
\end{definition}

Compare this with the pointwise definition, where $\delta$ may vary with $c$.
The key theorem below shows compactness is exactly what's needed to upgrade
continuity to uniform continuity.

\begin{theorem}[Heine--Cantor]
If $f: K \to \R$ is continuous and $K$ is compact, then $f$ is uniformly
continuous on $K$.
\end{theorem}

A quick numerical check --- verifying the $\delta/\eps$ bound on a sample function:

\begin{lstlisting}
def is_uniformly_cont(f, eps, xs):
    for delta in grid:
        if all(abs(f(x) - f(y)) < eps
               for x, y in pairs(xs, delta)):
            return delta
\end{lstlisting}

On the grid tested here, the search returns $\delta = 0.02$, well within the
tolerance we assumed at the start of the lecture. Note that this value depends
on the grid resolution, not on any particular point $c$ --- consistent with
the definition above.

\end{document}
